\chapter{Risultati e discussione}
\label{chap:risultati-discussione}

% ==================================================
% 7.1 Risultati ottenuti
% ==================================================

\section{Risultati ottenuti}
\label{sec:risultati-ottenuti}

% Sintetizzare solo risultati supportati dal prototipo e dalla
% validazione effettivamente eseguita.
%
% Non inventare metriche, percentuali, tempi di esecuzione,
% risultati sperimentali o feedback utenti non raccolti.

% ==================================================
% 7.2 Discussione
% ==================================================

\section{Discussione delle scelte progettuali}
\label{sec:discussione-scelte}

% Discutere i trade-off principali:
% - precisione geografica e privacy;
% - modularità e complessità del monorepo;
% - provider esterni e fallback;
% - catalogazione assistita e controllo umano;
% - prototipo universitario e limiti produttivi.

% ==================================================
% 7.3 Limiti
% ==================================================

\section{Limiti del prototipo}
\label{sec:limiti-prototipo}

% Distinguere limiti implementativi, metodologici e di
% validazione.
%
% Indicare come limite ogni funzionalità non implementata o
% non validata, senza presentarla come risultato.

% ==================================================
% 7.4 Sviluppi futuri
% ==================================================

\section{Sviluppi futuri}
\label{sec:sviluppi-futuri}

% Gli sviluppi futuri devono essere coerenti con lo stato
% reale del repository. Se una funzionalità viene implementata
% in futuro, rimuoverla da questa sezione e documentarla nel
% capitolo appropriato.
